% Bagean: metodhe
% Bab: induksi
% Pérangan: dhéfinisi-induktif

\documentclass[../../../include/open-logic-section]{subfiles}

\begin{document}

\olfileid{mth}{ind}{idf}

\olsection{Dhéfinisi Induktif}

Ing logika kita kerep ndhéfinisikaké jinis
obyèk \emph{kanthi induktif}: kita nemtokaké
aturan kang njlèntrèhaké obyèk apa kang
kalebu jinis iki lan kepriyé obyèk anyar
saka jinis iki dibangun saka obyèk lawas
kang padha jinisé. Upamané, rerangken
simbol tartamtu, kaya term lan !!{formula}s
ing sawijining basa, kerep didhéfinisikaké
kanthi induksi. Minangka conto prasaja,
delengen rerangken kang kasusun saka aksara
$\mathrm{a}$, $\mathrm{b}$, $\mathrm{c}$,
$\mathrm{d}$, simbol~$\circ$, lan kurung
siku $[$ lan $]$, kaya
``$[[\mathrm{c} \circ \mathrm{d}][$'',
``$[\mathrm{a}[]\circ]$'',
``$\mathrm{a}$'' utawa
``$[[\mathrm{a} \circ \mathrm{b}]\circ
\mathrm{d}]$''. Mbokmenawa panjenengan
rumangsa ana kang ``salah'' ing rong
rerangken kapisan: ing kang kapisan,
kurung sikuné ora ``imbang'' babar pisan;
panjenengan uga bisa rumangsa yèn
``$\circ$'' kuduné ``nyambungaké''
ekspresi kang dhéwé nduwèni teges.
Rerangken katelu lan kapapat katon
luwih trep: saben ``$[$'' nduwèni
pasangan nutup ``$]$'' yèn kurung
iku ana, lan ing saben $\circ$
ana ekspresi kang ``trep'' ing
loro sisihé, dikurung sapasang kurung.

Kita kepéngin nemtokaké kanthi persis
apa kang kalebu ``term trep''.
Kaping pisan, saben aksara kang madeg
dhéwé iku trep. Saliyane aksara tunggal,
term trep kudu awujud ``$[t \circ s]$'',
kanthi $s$ lan $t$ dhéwé trep.
Kosok baliné, yèn $t$ lan $s$
trep, kita bisa nggawe term trep
anyar kanthi nyeleh $\circ$
ing antarané lan ngubengi nganggo
sapasang kurung siku. Operasi iki
bisa digunakaké kanggo
\emph{ndhéfinisikaké} himpunan
term trep. Iki \emph{dhéfinisi
  induktif}.

\begin{defn}[Term trep]
  Himpunan \emph{term trep}
  didhéfinisikaké kanthi induktif:
  \begin{enumerate}
  \item Saben aksara $\mathrm{a}$,
    $\mathrm{b}$, $\mathrm{c}$,
    $\mathrm{d}$ iku term trep.
  \item Yèn $s_1$ lan $s_2$
    term trep, mula $[s_1 \circ s_2]$
    uga term trep.
  \item Ora ana liya kang kalebu
    term trep.
  \end{enumerate}
\end{defn}

Dhéfinisi iki nyatakaké yèn sawijining
obyèk kalebu term trep yèn lan mung
yèn bisa dibangun miturut syarat
(1) lan~(2) liwat cacah langkah winates.
Ing langkah kapisan, kita mbangun
kabèh term trep kang mung aksara
tunggal:
\[
\mathrm{a}, \mathrm{b}, \mathrm{c}, \mathrm{d}
\]
Ing langkah kapindho, kita nrapaké
(2) marang term kang wis dibangun.
Kita olèh
\[
[\mathrm{a} \circ \mathrm{a}], [\mathrm{a} \circ \mathrm{b}],
[\mathrm{b} \circ \mathrm{a}], \dots, [\mathrm{d} \circ \mathrm{d}]
\]
kanggo kabèh pasangan aksara.
Ing langkah katelu, kita nrapaké
(2) maneh marang pasangan term
trep kang wis dibangun. Kita
olèh term trep anyar kaya
$[\mathrm{a} \circ [\mathrm{a} \circ
    \mathrm{a}]]$, ing kéné
$t$ yaiku $\mathrm{a}$ saka
langkah~1 lan $s$ yaiku
$[\mathrm{a} \circ \mathrm{a}]$
saka langkah~2; lan
$[[\mathrm{b} \circ
    \mathrm{c}] \circ [\mathrm{d} \circ \mathrm{b}]]$
kang dibangun saka rong term
$[\mathrm{b} \circ \mathrm{c}]$
lan $[\mathrm{d}
  \circ \mathrm{b}]$ saka langkah~2.
Saterusé uga mangkono.
Klausa (3) nyegah obyèk kang
ora dibangun kanthi cara iki
melu mlebu himpunan term trep.

Élinga yèn kita durung mbuktekaké
kabèh rerangken simbol kang ``katon''
trep pancèn trep miturut dhéfinisi iki.
Nanging kuduné cetha yèn kabèh
kang bisa kita bangun pancèn
``katon trep'': kurung sikuné
imbang, lan $\circ$ nyambungaké
pérangan kang dhéwé trep.

Ciri pokok dhéfinisi induktif yaiku
yèn panjenengan arep mbuktekaké
pratelan tumrap kabèh term trep,
dhéfinisi iku nuduhaké kasus endi
kang kudu ditliti. Upamané, yèn
$t$ diwènèhaké minangka term trep,
dhéfinisi induktif nuduhaké
wujud kang mungkin kanggo $t$:
$t$ bisa awujud aksara, utawa
$[s_1 \circ s_2]$ kanggo pasangan
term trep $s_1$ lan~$s_2$.
Amarga klausa (3), mung iku
kamungkinané.

Nalika mbuktekaké pratelan
tumrap kabèh anggota himpunan
kang didhéfinisikaké kanthi induktif,
induksi kuwat mligi wigati.
Upamané, kita arep mbuktekaké
yèn ing saben term trep kanthi
dawa~$n$, cacah $[$ iku~$< n/2$.
Iki bisa diwaca minangka pratelan
kanggo kabèh~$n$: kanggo saben $n$,
cacah $[$ ing sembarang term trep
kanthi dawa~$n$ iku $< n/2$.

\begin{prop}
  Kanggo sembarang $n$, cacah $[$
  ing term trep kanthi dawa~$n$
  iku $< n/2$.
\end{prop}

\begin{proof}
Kanggo mbuktekaké kanthi induksi
(kuwat), kita kudu nuduhaké yèn
pratelan kondisional iki bener:
\begin{quote}
  Yèn kanggo saben $l < k$,
  saben term trep kanthi dawa~$l$
  nduwèni $< l/2$ simbol $[$,
  mula saben term trep kanthi
  dawa~$k$ nduwèni $< k/2$
  simbol $[$.
\end{quote}
Kanggo mbuktekaké kondisional
iki, anggepen antesedèné bener:
kanggo sembarang $l<k$,
term trep kanthi dawa~$l$
ngemot $< l/2$ simbol $[$.
Anggepan iki diarani
hipotesis induksi. Kita kudu
nuduhaké bab kang padha kanggo
term trep kanthi dawa~$k$.

Anggepen $t$ term trep kanthi
dawa~$k$. Amarga term trep
didhéfinisikaké kanthi induktif,
ana rong kasus: (1)~$t$ mung
aksara, utawa (2)~$t$ awujud
$[s_1 \circ s_2]$ kanggo
term trep $s_1$ lan~$s_2$.
\begin{enumerate}
\item Yèn $t$ aksara, mula
  $k = 1$, lan cacah $[$
  ing $t$ yaiku~$0$.
  Amarga $0 < 1/2$, pratelan
  iki bener.
\item Yèn $t$ awujud
  $[s_1 \circ s_2]$ kanggo
  term trep $s_1$ lan~$s_2$,
  tetepna $l_1$ minangka
  dawa~$s_1$ lan $l_2$
  minangka dawa~$s_2$.
  Mula dawa~$k$ saka $t$
  yaiku $l_1+l_2+3$:
  dawa $s_1$ lan $s_2$
  ditambah telung simbol
  $[$, $\circ$, $]$.
  Amarga $l_1+l_2+3$
  tansah luwih gedhé tinimbang
  $l_1$, kita olèh $l_1 < k$.
  Semono uga $l_2
  < k$. Dadi hipotesis induksi
  bisa ditrapaké marang
  $s_1$ lan~$s_2$:
  cacah~$m_1$ simbol $[$
  ing $s_1$ iku $< l_1/2$,
  lan cacah~$m_2$ simbol $[$
  ing~$s_2$ iku $< l_2/2$.

  Cacah $[$ ing $t$
  yaiku cacah $[$ ing~$s_1$,
  ditambah cacah $[$ ing~$s_2$,
  ditambah~$1$, yaiku
  $m_1 + m_2 + 1$.
  Amarga $m_1
  < l_1/2$ lan $m_2 < l_2/2$,
  kita olèh:
  \[
  m_1 + m_2 + 1 < \frac{l_1}{2} + \frac{l_2}{2} + 1 = \frac{l_1+l_2+2}{2} < \frac{l_1+l_2+3}{2} = k/2.
  \]
\end{enumerate}
Ing saben kasus, cacah $[$
ing $t$ kabuktèn $< k/2$
kanthi hipotesis induksi.
Miturut induksi kuwat, proposisi
iku katut.
\end{proof}

\begin{prob}
  Dhéfinisikna himpunan term
  supertrep mangkéné:
  \begin{enumerate}
  \item Saben aksara
    $\mathrm{a}$, $\mathrm{b}$,
    $\mathrm{c}$, $\mathrm{d}$
    iku term supertrep.
  \item Yèn $s$ term supertrep,
    mula $[s]$ uga mangkono.
  \item Yèn $s_1$ lan $s_2$
    term supertrep, mula
    $[s_1 \circ s_2]$ uga mangkono.
  \item Ora ana liya kang
    kalebu term supertrep.
  \end{enumerate}
  Tuduhna yèn cacah $[$
  ing term supertrep~$t$
  kanthi dawa~$n$ iku
  $\le n/2 +1$.
\end{prob}

\end{document}
